\section{Constructing the measure-precompact Daugavet space}
\label{kw:section}

The preceding results apply to a space with two rather different
properties: its unit ball is small for convergence in measure, while
its slices contain vectors almost at distance two from any prescribed
unit vector.
We now construct such a space.  The construction is due to
Kadets and Werner \cite[Lemma~2.4 and Theorem~2.5]{KW2004}, who refine
the method of Bourgain and Rosenthal \cite{BourgainRosenthal1980}.
We give all the estimates needed for the finite-dimensional enlargement
and its iteration. The explicit multipliers below and the constants
$10$, $54$, and $108$ belong to the quantitative realization given here;
the enlargement and iteration scheme is the credited construction.

The idea is to multiply an old unit vector by a positive random variable
which has mean one but is small with high probability.  Each product has
small $L^1$ overlap with every old unit vector, so the norm of their sum is
almost two.  Averaging sufficiently many products recovers the old vector
in norm.  The main additional task is to control every vector in the
new unit ball, which permits the iteration to preserve precompactness
in measure.

Throughout this section,
\[
 (\Omega,\PP)=([0,1]^{\N},\lambda^{\N}),
 \qquad L^1=L^1(\Omega,\PP;\R),
\]
and $\norm{\cdot}$ denotes the $L^1$ norm.  Write $B(G)$ and $S(G)$ for
the closed unit ball and unit sphere of a subspace $G$.  A space depends
on finitely many coordinates if one common finite set of coordinates
determines all its elements.

Recall the metric from \eqref{eq:measure-metric}, on measurable functions
modulo equality almost everywhere:
\begin{equation}\label{kw:metric}
 d(f,g)=\inf\{a>0:\PP(\abs{f-g}>a)<a\}.
\end{equation}
This metric induces convergence in measure.  For example, if
$\PP(\abs{f-g}>a)<a$ and $\PP(\abs{g-h}>b)<b$, the union bound gives
$\PP(\abs{f-h}>a+b)<a+b$; taking infima proves the triangle inequality.
The identity $d(f,g)=0$ is equivalent to $f=g$ almost everywhere, and
the convergence assertion follows directly from the thresholds in
\eqref{kw:metric}.  Markov's inequality also gives
$d(f,g)\le\norm{f-g}^{1/2}$, so norm convergence implies convergence
for $d$.

\begin{theorem}[Kadets--Werner construction]\label{kw:existence}
There is an infinite-dimensional closed real subspace $E\subseteq L^1$
such that $B(E)$ is totally bounded for $d$ and has the following slice
property: for every $w\in S(E)$, every nonempty slice
\[
 H=\{y\in B(E):\ell(y)>\alpha\},\qquad \ell\in E^*,\quad\alpha\in\R,
\]
and every $\rho>0$, there is $y\in H\cap S(E)$ such that
$\norm{w+y}\ge2-\rho$.  In particular, $E$ has the Daugavet property.
\end{theorem}

\subsection{The finite-dimensional enlargement}

We need the new unit vectors to be almost additive in norm with every
old unit vector, and their averages to approximate selected old vectors.
The third conclusion below keeps the entire new unit ball close, in
measure, to the old unit ball.

\begin{lemma}\label{kw:enlargement}
Suppose $G\ne\{0\}$ is a finite-dimensional subspace of $L^1$ depending
on finitely many coordinates, $m\ge1$, $u_1,\dots,u_m\in S(G)$, and $0<h<1$.
There are a finite-dimensional $G'\supseteq G$ depending on finitely
many coordinates, an integer $n\ge1$, and vectors
$v_{kj}\in S(G')$ $(1\le k\le m,\ 1\le j\le n)$ such that
\begin{align}
 \norm{w+v_{kj}}&\ge2-h
 &&(w\in S(G),\ k\le m,\ j\le n),\label{kw:enlarge-a}\\
 \norm{u_k-\frac1n\sum_{j=1}^n v_{kj}}&\le h
 &&(k\le m),\label{kw:enlarge-b}\\
 \text{for every }g\in B(G')\text{ there is }\psi_0\in B(G)
 &\text{ with }d(g,\psi_0)\le h.\label{kw:enlarge-c}
\end{align}
\end{lemma}

To prove the lemma, we will use independent positive multipliers whose
mean is one and whose mass is concentrated near zero.  The next estimate
controls their linear combinations; this is what the whole-ball
conclusion requires beyond control of the selected vectors themselves.

\subsection{A heavy-tail estimate}

The following variable has mean one for every $p>1$, although it converges
to zero in probability as $p\downarrow1$.  The estimate for differences
will control arbitrary linear combinations in an enlarged unit ball.

\begin{lemma}\label{kw:heavy-tail}
Let $1<p\le3/2$, $s=(p-1)/p$, and $f=sV^{-1/p}$, where $V$ is uniform
on $(0,1)$.  Let $Y_1,\dots,Y_n$ be independent variables, each distributed
as the difference of two independent copies of $f$.  For real scalars
$a_1,\dots,a_n$, put
$A=(\sum_{j=1}^n\abs{a_j}^{p})^{1/p}$ and $Z=\sum_{j=1}^n a_jY_j$.
Then $\E f=1$ and
\begin{equation}\label{kw:heavy-bounds}
 \PP(\abs Z>u)\le10(sA/u)^p\quad(u>0),
 \qquad \E\abs Z\ge A/54.
\end{equation}
\end{lemma}

The lower estimate bounds $A$ from an $L^1$ bound on $Z$;
the upper estimate then gives concentration in probability uniformly in $n$.

\begin{proof}
Direct integration gives $\E f=sp/(p-1)=1$, and
$\PP(f>u)=(s/u)^p$ for $u\ge s$.  For all $u>0$ the same expression
is an upper bound.  If $Y=f-f'$ for independent copies, then $Y$ is
symmetric and $\PP(\abs Y>u)\le2(s/u)^p$.
For fixed $u>0$, the variables
$X_j=a_jY_j\mathbf1_{\{\abs{a_jY_j}\le u\}}$ are independent and
centered, and
\[
 \E X_j^2
 \le\int_0^u 2v\,\PP(\abs{a_jY_j}>v)\,dv
 \le\frac{4(s\abs{a_j})^p u^{2-p}}{2-p}.
\]
Except on the event that some $\abs{a_jY_j}>u$, the sums $Z$ and
$\sum_jX_j$ agree.  The union bound and Chebyshev's inequality therefore
give
\[
 \PP(\abs Z>u)
 \le\left(2+\frac4{2-p}\right)(sA/u)^p
 \le10(sA/u)^p.
\]

For the lower bound, when $u\ge s$, the event
$\{f>u+2s,\ f'\le2s\}$ gives
\[
 \PP(\abs Y>u)
 \ge\left(\frac{s}{u+2s}\right)^p(1-2^{-p})
 \ge\frac1{18}(s/u)^p.
\]
Here $u+2s\le3u$, $1-2^{-p}\ge1/2$, and $3^p\le9$.
Suppose $A>0$ and omit zero coefficients.  Since $\abs{a_j}\le A$,
independence implies, for $u\ge sA$,
\begin{align*}
 \PP\left(\max_j\abs{a_jY_j}>u\right)
 &\ge1-\exp\left(-\frac1{18}(sA/u)^p\right)\\
 &\ge\frac1{36}(sA/u)^p.
\end{align*}
The last inequality uses $1-e^{-v}\ge v/2$ for $0\le v\le1$.
Adjoin independent signs $\varepsilon_j\in\{-1,1\}$ of equal probability.
For each fixed realization of the terms, averaging first over the sign
of a term of largest absolute value gives
\[
 \E_{\varepsilon}\abs{\sum_j\varepsilon_j a_jY_j}
 \ge\max_j\abs{a_jY_j},
\]
because $(\abs{b+c}+\abs{b-c})/2\ge\abs c$.
The signed sum has the same law as $Z$, by independence and symmetry.
Consequently,
\[
 \E\abs Z\ge\E\max_j\abs{a_jY_j}
 \ge\frac1{36}\int_{sA}^\infty(sA/u)^p\,du
 =\frac{A}{36p}\ge\frac A{54}.
\]
If $A=0$, both assertions are immediate.
\end{proof}

\subsection{Proof of the enlargement}

The scalar estimates now let us control all linear combinations in the
new space.  We prove the three conclusions of Lemma~\ref{kw:enlargement}
in the order that permits the number of multipliers to be chosen last.

\begin{proof}[Proof of Lemma~\ref{kw:enlargement}]
We first fix the distribution of the new multipliers and prove the first
and third conclusions uniformly in their number.  Only then will we
choose the number of multipliers to obtain the averaging conclusion.

\emph{Parameter choices.}
Choose $c\in(0,h/4)$, and then $r,q>0$ sufficiently small that
\begin{equation}\label{kw:rq-choice}
 mq<1,\qquad
 mq+h^{-1}\bigl(2mr+(1-mq)^{-1}-1\bigr)<h.
\end{equation}
With $s=(p-1)/p$, choose $1<p\le3/2$ close enough to $1$ that
\begin{equation}\label{kw:p-choice}
 s<c,\qquad 2\bigl((s/c)^p+c\bigr)<h,
 \qquad 10(108s/r)^p<q.
\end{equation}
These choices are possible because $s\to0$ as $p\downarrow1$.
We fix $p$ now; the integer $n$ will be chosen last.

\emph{Near-additivity with the old space.}
For any prospective $n$, take $mn$ independent copies $f_{kj}$ of the
variable $f$ in Lemma~\ref{kw:heavy-tail}, each on a fresh coordinate,
and put
\[
 v_{kj}=u_kf_{kj},\qquad
 G'_n=\operatorname{span}\bigl(G\cup\{v_{kj}:k\le m,\ j\le n\}\bigr).
\]
Independence from the old coordinates gives $\norm{v_{kj}}=1$.
For every $w\in S(G)$,
\begin{align*}
 \norm{w+v_{kj}}
 &\ge2-2\E\min(\abs w,\abs{u_k}f_{kj})\\
 &\ge2-2\bigl(\PP(f>c)+c\bigr)>2-h.
\end{align*}
Indeed, on $\{f_{kj}>c\}$ bound the minimum by $\abs w$, and on its
complement by $c\abs{u_k}$; independence then gives the last bound.
Thus \eqref{kw:enlarge-a} holds for every $n$.

\emph{Approximation of the whole unit ball.}
We next prove \eqref{kw:enlarge-c}, also uniformly in $n$.
Fix $g\in B(G'_n)$ and choose any representation
\[
 g=x+\sum_{k=1}^m u_kz_k,\qquad
 x\in G,\quad z_k=\sum_{j=1}^n a_{kj}f_{kj}.
\]
Fix $k$.  On an auxiliary product probability space resample all the
variables in the $k$th group, leaving the old coordinates and all other
groups unchanged.  Write $z'_k$ for the resampled sum and $g^{(k)}$
for the resulting function.  The new variables are independent copies,
so $g^{(k)}$ has the same law as $g$, and
\[
 \E\abs{z_k-z'_k}
 =\E\abs{u_k(z_k-z'_k)}
 =\E\abs{g-g^{(k)}}\le2.
\]
The first equality uses $\norm{u_k}=1$ and independence of the old
coordinates.  Lemma~\ref{kw:heavy-tail}, applied to the differences of
the paired copies, therefore gives
\begin{equation}\label{kw:coefficient-bound}
 A_k=\left(\sum_j\abs{a_{kj}}^p\right)^{1/p}\le108,
 \qquad
 \PP(\abs{z_k-z'_k}>r)\le10(108s/r)^p<q.
\end{equation}
In particular, the bound does not depend on $n$ or on the chosen
representation of $g$.

The estimate for $z_k-z'_k$ says that $z_k$ is concentrated near some
constant.  Indeed, averaging the conditional probability over $z'_k$
gives a deterministic $c_k\in\R$ with
$\PP(\abs{z_k-c_k}>r)<q$.
These constants may depend on $g$; no simultaneous choice is needed.
Set
\[
 \psi=x+\sum_{k=1}^m c_ku_k\in G,
 \qquad D=\bigcap_{k=1}^m\{\abs{z_k-c_k}\le r\}.
\]
The event $D$ depends only on the fresh coordinates, has probability
at least $1-mq$, and on $D$ we have
$\abs{g-\psi}\le r\sum_k\abs{u_k}$.
Its independence from the old coordinates gives
\[
 1\ge\E(\abs g\mathbf1_D)
 \ge\PP(D)(\norm\psi-mr),
 \qquad \norm\psi\le mr+(1-mq)^{-1}.
\]
At this point $\psi$ need not belong to the old unit ball.  Set
$\psi_0=\psi/\max\{1,\norm\psi\}\in B(G)$.  The cost of this normalization
is bounded by
\[
 \norm{\psi-\psi_0}
 \le mr+(1-mq)^{-1}-1.
\]
Markov's inequality restricted to $D$ and the bound on its complement
now yield
\begin{align*}
 \PP(\abs{g-\psi_0}>h)
 &\le mq+h^{-1}\bigl(mr+\norm{\psi-\psi_0}\bigr)\\
 &\le mq+h^{-1}\bigl(2mr+(1-mq)^{-1}-1\bigr)<h.
\end{align*}
This proves \eqref{kw:enlarge-c} for every $n$.

We have now proved both near-additivity and whole-ball approximation
without placing an upper bound on $n$.

\emph{Recovering the old vectors by averages.}
It remains to choose $n$ for \eqref{kw:enlarge-b}.
For independent copies $f_1,\dots,f_n$ of the fixed integrable variable
$f$, write $f^M=\min\{f,M\}$.  Truncating, centering, and applying the
variance bound to the bounded independent copies gives
\begin{equation}\label{kw:lln}
 \E\abs{\frac1n\sum_{j=1}^nf_j-1}
 \le2\E(f-f^M)+\sqrt{\operatorname{Var}(f^M)/n}.
\end{equation}
The first term tends to zero as $M\to\infty$; after $M$ is fixed the
second tends to zero as $n\to\infty$.  Choose $n$ so that the left side
is at most $h$.  All groups have the same law, and independence gives
\[
 \norm{u_k-\frac1n\sum_jv_{kj}}
 =\E\abs{1-\frac1n\sum_jf_{kj}}\le h
 \qquad(k\le m).
\]
Taking $G'=G'_n$ completes the proof.  This space uses just finitely many
additional coordinates; the earlier two conclusions persist because
their estimates held for every finite $n$.
\end{proof}

\subsection{Passage to the closed space}

We apply the enlargement to finer and finer nets of the old unit sphere.
The summable errors keep the eventual unit ball precompact in measure;
the approximation by averages ensures that every slice eventually
contains one of the new unit vectors.

\begin{proof}[Proof of Theorem~\ref{kw:existence}]
Start with $G_1=\operatorname{span}\{1\}$ and put $h_b=2^{-b}$.
Given $G_b$, choose a finite $h_b$-net
$\{u_{b,k}:1\le k\le m_b\}$ of $S(G_b)$ in norm and apply
Lemma~\ref{kw:enlargement} with $h=h_b$.  Denote the resulting space by
$G_{b+1}$ and its new unit vectors by $v_{b,kj}$, $1\le j\le n_b$.
All stages depend on finitely many coordinates, so the construction
continues on $\Omega$.  Set
\[
 E=\overline{\bigcup_{b\ge1}G_b}^{\,\norm{\cdot}}\subseteq L^1.
\]
Each inclusion $G_b\subseteq G_{b+1}$ is strict: if any $v_{b,kj}$
belonged to $G_b$, substituting $w=-v_{b,kj}$ in
\eqref{kw:enlarge-a} would give $0\ge2-h_b$.
Thus $E$ is infinite-dimensional.

\emph{The entire unit ball is totally bounded in measure.}
The union
$\bigcup_bB(G_b)$ is norm dense in $B(E)$: if $x_j$ in the algebraic
union tends to $x\in B(E)$, then
$x_j/\max\{1,\norm{x_j}\}$ also tends to $x$.
Repeated use of \eqref{kw:enlarge-c} gives, for $a>b$ and $g\in B(G_a)$,
\[
 \inf_{\psi\in B(G_b)}d(g,\psi)
 \le\sum_{l=b}^{a-1}h_l
 \le\sum_{l=b}^\infty h_l=2^{1-b}.
\]
The bound is automatic for $a\le b$.  Distance to a fixed nonempty
set is a $1$-Lipschitz function for $d$, so norm density and
\eqref{kw:metric} extend it to every $g\in B(E)$.
The set $B(G_b)$ is norm compact and therefore $d$-compact.
Given $\eps>0$, choose $b$ with $2^{1-b}<\eps/2$ and a finite
$\eps/2$-net of $B(G_b)$ for $d$.  The preceding distance estimate
shows that this is an $\eps$-net of $B(E)$.

\emph{The Daugavet slice property.}
Now fix $w,H,\ell,\alpha,\rho$ as in Theorem~\ref{kw:existence}.
There is $u\in S(E)$ with $\ell(u)>\alpha$: the supremum of a linear
functional over the ball equals its supremum over the sphere, including
when the functional is zero.  Normalizing approximants from the
increasing union gives $u^{(b)},w^{(b)}\in S(G_b)$ with
$u^{(b)}\to u$ and $w^{(b)}\to w$ in norm.
Choose a net point $u_{b,k_b}$ with
$\norm{u_{b,k_b}-u^{(b)}}\le h_b$ and set
\[
 a_b=\frac1{n_b}\sum_{j=1}^{n_b}v_{b,k_bj}.
\]
By \eqref{kw:enlarge-b},
$\norm{a_b-u}\le2h_b+\norm{u^{(b)}-u}\to0$.
Thus $\ell(a_b)>\alpha$ for all sufficiently large $b$, and for each
such $b$ at least one summand $v_{b,k_bj}$ lies in $H$.
For every summand, \eqref{kw:enlarge-a} also gives
\[
 \norm{w+v_{b,k_bj}}
 \ge2-h_b-\norm{w-w^{(b)}}.
\]
Take $b$ large enough that the last two error terms sum to less than
$\rho$.  The summand in $H$ is the required unit vector.

Finally, Lemma~\ref{lem:daugavet-slices} turns this slice property into
$\norm{I+T}=1+\norm T$ for every rank-one operator $T$ on $E$.
\end{proof}
